\section{Reference Conditions}
\label{SEC:Reference_Conditions}

The user-specified reference conditions are used to initialize
the volume solution and to compute the non-dimensional variables
used within the flow solver.
%
The reference conditions typically represent the freestream flow.
%
These conditions include the gas properties:
ratio of specific heats ($\gamma$),
gas constant ($R_{gas}$),
Prandtl number ($Pr$),
Schmidt number ($Sc$), and
coefficients to the Sutherland viscosity law.
%
They also include the flow conditions:
Mach number, Reynolds number, static pressure, temperature, density,
as well as
\key*[key:AngleOfAttack]{angle of attack}[flow angles!angle of attack]
$\alpha$ in the $xy$-plane and
\key*[key:AngleOfSideslip]{angle of sideslip}[flow angles!angle of sideslip]
$\beta$ in the $xz$-plane.
%
%%% Index aliases:
\index{angle of attack|see{flow angles}}
\index{angle of sideslip|see{flow angles}}
%
These flow angles are depicted in \fref{FIG:FlowAngle} and used
to compute the velocity vector components as follows:
\begin{align}
  u =& U \cos{(\beta)} \cos{(\alpha)} \label{EQN:FlowAngleVelX} \\
  v =& U \cos{(\beta)} \sin{(\alpha)} \label{EQN:FlowAngleVelY} \\
  w =& U \sin{(\beta)}                \label{EQN:FlowAngleVelZ}
\end{align}
%
For positive angle of attack $\alpha$ between 0 and 90, the velocity will
have a positive component in the $y$-direction.
%
For positive angle of sideslip $\beta$ between 0 and 90, the velocity will
have a positive component in the $z$-direction.

%%% Replace the tikz image with tecplot image.
%
%\newpage
%
%\begin{figure}[htb]
%  \begin{centering}
%    % trim {l,b,r,t}
%    \insertimage[trim={0.5in 0.90in 0.0in 0.35in},clip,width=6.0in]{tikz/tikz_aoa_cfd}
%    \caption{Flow angles.}
%    \label{FIG:FlowAngle}
%  \end{centering}
%\end{figure}

\begin{figure}[htb]
  \begin{centering}
    % trim {l,b,r,t}
    \insertimage[width=6.0in]{aoa/BELL_X1/aoa_vectors}
    \caption{Flow angles.}
    \label{FIG:FlowAngle}
  \end{centering}
\end{figure}

Reference conditions are specified in the \inp[nmlReference]{nml\_reference}
input namelist block.
%
Details regarding the velocity and length scale used to compute the
Reynolds number can be found in the description of the
\inp[nmlReference::Reynolds]{Reynolds} input variable.

\newcommand{\refval}[1]{\ensuremath{#1_\infty}}
\newcommand{\dimval}[1]{\ensuremath{#1}}
\newcommand{\nonval}[1]{\ensuremath{#1^*}}
\newcommand{\taustar}{\ensuremath{\nonval{\scriptstyle \mathcal{T}}}}
%\newcommand{\myfrac}[2]{\ensuremath{\frac{#1}{#2}}}
\newcommand{\myfrac}[2]{\ensuremath{{#1}/{#2}}}
Internally, GFR assumes
$\refval{L}=1.0\unit{m}$,
$\refval{t}=\refval{L}/\refval{a}$,
and the following non-dimensionalization:
%
\begin{align*}
%   \nonval{\mathfrak{t}}    =& \myfrac{\dimval{t   }}{\refval{t   }}
%&  \nonval{\rho} =& \myfrac{\dimval{\rho}}{\refval{\rho}}
%&  \nonval{e_t}  =& \myfrac{\dimval{e_t }}{\refval{a   }^2}
%\\
%   \nonval{{\scriptstyle \mathbb{T}}}    =& \myfrac{\dimval{t   }}{\refval{t   }}
%&  \nonval{\rho} =& \myfrac{\dimval{\rho}}{\refval{\rho}}
%&  \nonval{e_t}  =& \myfrac{\dimval{e_t }}{\refval{a   }^2}
%\\
    \taustar      =& \myfrac{\dimval{t   }}{\refval{t   }}
 &  \nonval{\rho} =& \myfrac{\dimval{\rho}}{\refval{\rho}}
 &  \nonval{e_t}  =& \myfrac{\dimval{e_t }}{\refval{a   }^2}
 \\
%   \nonval{{\scriptscriptstyle \mathcal{T}}}    =& \myfrac{\dimval{t   }}{\refval{t   }}
%&  \nonval{\rho} =& \myfrac{\dimval{\rho}}{\refval{\rho}}
%&  \nonval{e_t}  =& \myfrac{\dimval{e_t }}{\refval{a   }^2}
%\\
%   \nonval{\tau}    =& \myfrac{\dimval{t   }}{\refval{t   }}
%&  \nonval{\rho} =& \myfrac{\dimval{\rho}}{\refval{\rho}}
%&  \nonval{e_t}  =& \myfrac{\dimval{e_t }}{\refval{a   }^2}
%\\
%   \nonval{t}    =& \myfrac{\dimval{t   }}{\refval{t   }}
%&  \nonval{\rho} =& \myfrac{\dimval{\rho}}{\refval{\rho}}
%&  \nonval{e_t}  =& \myfrac{\dimval{e_t }}{\refval{a   }^2}
%\\
    \nonval{x}    =& \myfrac{\dimval{x   }}{\refval{L   }}
 &  \nonval{p}    =& \myfrac{\dimval{p   }}{\refval{\rho}\refval{a}^2}
 &  \nonval{\mu}  =& \myfrac{\dimval{\mu }}{\refval{\mu }}
 \\
    \nonval{u}    =& \myfrac{\dimval{u   }}{\refval{a   }}
 &  \nonval{T}    =& \myfrac{\dimval{T   }}{\refval{T   }}
\end{align*}

%%%%%%%%%%%%%%%%%%%%%%%%%%%%%%%%%%%%%%%%%%%%%%%%%%%%%%%%%%%%%%%%%%%%%%%%%%%%%%%%

\noindent\NOTE
 
The nondimensional time, \taustar, is only used in two places within GFR: the
second column of the iteration data written to the GFR output and the input
\inp[nmlTimestep::Timestep]{Timestep} are both nondimensionalized the same as
\taustar.
 
There is another nondimensional time that is used more frequently in GFR.
%
This nondimensional time, \nonval{t}, uses the input
\inp[nmlOutput::TimeScalingFactor]{Time\_Scaling\_Factor} to scale the
simulation time by something that may be more useful for a given problem,
\forexample flow-through-time for a channel flow or nozzle-diameter-time for a
jet.
%
These two nondimensional times are related through the equation
\begin{equation}
 %\nonval{t} = \,{\taustar}\cdot\refval{t}\cdot[\inptext{Time\_Scaling\_Factor}]
  \nonval{t} = {\taustar}\cdot\refval{t}\cdot[\inptext{Time\_Scaling\_Factor}]
             = \dimval{t}\cdot[\inptext{Time\_Scaling\_Factor}]
  \label{EQN:taustar_to_tstar}
\end{equation}
%
References to the term \quoterm{tstar} in the inputs, output, source code,
documentation, and elsewhere in the GFR universe are referring to the value
\nonval{t}.
